\chapter{The North American Map}\label{ch:nw:the-north-american-map}

Three data centres in New Jersey, a few tens of kilometres apart, hold the matching engines of the NYSE group, of Nasdaq and of Cboe's U.S.
equities and options exchanges; a fourth, in Aurora, west of Chicago, holds CME's Globex. Light in a vacuum needs \qty{3.94}{\milli\second}
to cross from Aurora to Nasdaq's building in Carteret, light in glass \qty{5.76}{\milli\second}. In 2010 the fastest fibre between
Chicago and Carteret was quoted at \qty{6.65}{\milli\second}; in 2016 a microwave network delivered Aurora's data to Carteret in
\qty{3.982}{\milli\second}, 1.1\% above the time light needs in air along the shortest path on the Earth's surface. The history of North
American trading networks is the story of closing that gap, and this chapter measures it.

Chapter~9 described what a venue's building sells. This chapter puts the buildings on a map: which venue matches where, how far apart the
buildings are, how long light needs between them in vacuum, in air and in glass, and how close to those floors the published routes come.
Every coordinate is a ledger row with its source, every distance is computed on the WGS-84 ellipsoid, and every map is drawn from
computed coordinates. Chapters~11 and~12 do the same for Europe and Asia, and chapters~13 and~14 look inside the routes.

\section{The New Jersey triangle}

\begin{definition}[Geodesic distance]\label{def:nw:the-north-american-map:geodesic}
The \emph{geodesic distance}\index{geodesic distance} between two points on the Earth is the length of the shortest path between them along
the surface of a reference ellipsoid (for maps and satellite positioning, WGS-84), computed from their latitudes and longitudes.
\end{definition}

The geodesic is the path a cable or a line of towers would follow if nothing were in the way: no rivers, no private land, no hills, no
permits. It is the natural yardstick for a trading route, and it is not the shortest path through space: the straight chord through the
Earth between Aurora and Carteret is \qty{1.68}{\kilo\metre} shorter, \qty{5.6}{\micro\second} of light, and nobody is going to dig it.

\begin{dated}{2026-09}{Who matches where in North America}\label{dat:nw:the-north-american-map:sites}
\begin{itemize}
\item \textbf{Mahwah, New Jersey} (1700 McArthur Boulevard): ICE's U.S. Liquidity Center, which hosts the NYSE group's markets and the
  consolidated feeds of the CTA and OPRA plans.
\item \textbf{Carteret, New Jersey} (NY11, 1400 Federal Boulevard, an Equinix building): Nasdaq's data centre, in expansion (NY11-4) for
  power and cabinets.
\item \textbf{Secaucus, New Jersey} (Equinix NY5, 800 Secaucus Road): Cboe's U.S. equities and options exchanges; Equinix NY4, at 755
  Secaucus Road, a few hundred metres away, is one of the points to which ICE delivers Toronto's market data.
\item \textbf{Aurora, Illinois} (2905 Diehl Road): CME Group's largest data centre, leased from CyrusOne until 2031; CME and Google Cloud
  are building a private cloud region and a co-location facility in Aurora.
\item \textbf{Chicago} (350 Cermak): Cboe's secondary data centre for all its U.S. platforms.
\item \textbf{Markham, Ontario}: TMX's main data centre (TSX, TSX Venture, TSX Alpha, Montreal Exchange), run by TMX; its street address
  could not be confirmed from a primary source and is flagged approximate.
\end{itemize}
\end{dated}

\Cref{fig:nw:the-north-american-map:nj} is the triangle, drawn from the coordinates of the site table. Mahwah is the northern corner, 55.6
kilometres from Carteret and 34.5 from Secaucus; Secaucus and Carteret are 25.9 kilometres apart. NY4 and NY5 stand 320 metres from each
other, one dot at this scale. In vacuum the triangle's sides take 185, 115 and \qty{86}{\micro\second}; in straight fibre, half as much again.

\begin{omfigure}
\begin{tikzpicture}
\begin{axis}[axis equal image,width=10cm,xmin=-30,xmax=58,ymin=-32,ymax=38,xlabel={km east of 74.15\textdegree{} W},
  ylabel={km north of 40.8\textdegree{} N},tick label style={font=\footnotesize},label style={font=\small},grid=major,grid style={gray!20},
  table/col sep=comma,unbounded coords=jump]
\addplot[omDef,thick,no markers] table[x=x,y=y]{figdata/networks/10-the-north-american-map/nj_links.csv};
\addplot[only marks,mark=*,mark size=2.6pt,omThm,visualization depends on={value \thisrow{anchor}\as\anc},nodes near coords,
  point meta=explicit symbolic,every node near coord/.append style={font=\footnotesize,anchor=\anc,black}]
  table[x=x,y=y,meta=name]{figdata/networks/10-the-north-american-map/nj_sites.csv};
\addplot[only marks,mark=none,nodes near coords,point meta=explicit symbolic,
  every node near coord/.append style={font=\scriptsize,fill=white,inner sep=1pt,anchor=center,text=omMeth}]
  table[x=x,y=y,meta=text]{figdata/networks/10-the-north-american-map/nj_labels.csv};
\end{axis}
\end{tikzpicture}

\omcaption{The New Jersey triangle from computed coordinates: each site from its sourced street address and a cited geocoder
(\cref{dat:nw:the-north-american-map:sites}), each side labelled with its WGS-84 \omterm{def:nw:the-north-american-map:geodesic}{geodesic distance} and light-time in vacuum. The projection is
a local equirectangular one centred on 40.8\textdegree{}~N, 74.15\textdegree{}~W. Data: \texttt{fig\_map.py} on
\texttt{firm.geomap}.}\label{fig:nw:the-north-american-map:nj}
\end{omfigure}

The triangle sets the arithmetic of U.S. equity and options trading. The consolidated feed of the CTA plan is produced in Mahwah, so a
quote change at Nasdaq in one of that plan's stocks must travel from Carteret to Mahwah, be processed, and travel back before a firm in
Carteret sees it in the consolidated feed: at least twice \qty{185.4}{\micro\second}, \qty{371}{\micro\second}, of light in vacuum, before any processing. The same firm reads
Nasdaq's direct feed in its own building. The gap between the direct feeds and the securities information processor, which One Quant
Book~1 describes, is partly geography, and no amount of engineering at the SIP removes the geographic part.

The triangle also sets the timing of an order split across venues. A router in Carteret that sends one piece to Cboe in Secaucus and one to
the NYSE in Mahwah reaches Secaucus first: \qty{126.6}{\micro\second} of straight fibre against \qty{271}{\micro\second}. Synchronised routing
delays the Secaucus piece by the difference, here about \qty{144}{\micro\second} along geodesics, in practice by the difference in measured
latencies of the firm's real routes, which are longer. Every latency race among the three buildings (latency arbitrage, in One Quant Book~11's
terms) is fought on these distances.

\begin{table}[ht]
\centering\small
\begin{tabular}{llrrrr}
\toprule
From & To & Geodesic & Vacuum & Air & Fibre \\
 & & (km) & (\si{\micro\second}) & (\si{\micro\second}) & (\si{\micro\second}) \\
\midrule
Mahwah & Carteret & 55.6 & 185.4 & 185.4 & 271.0 \\
Mahwah & Secaucus NY4 & 34.5 & 114.9 & 115.0 & 168.0 \\
Secaucus NY4 & Carteret & 25.9 & 86.3 & 86.3 & 126.2 \\
Aurora & Carteret & 1\,180.9 & 3\,939.0 & 3\,940.1 & 5\,758.8 \\
Aurora & Mahwah & 1\,179.0 & 3\,932.6 & 3\,933.8 & 5\,749.4 \\
Aurora & Secaucus NY5 & 1\,191.1 & 3\,973.1 & 3\,974.3 & 5\,808.7 \\
Aurora & 350 Cermak & 52.3 & 174.5 & 174.5 & 255.1 \\
350 Cermak & Carteret & 1\,129.2 & 3\,766.7 & 3\,767.8 & 5\,506.9 \\
Markham & Mahwah & 523.7 & 1\,747.0 & 1\,747.6 & 2\,554.2 \\
Markham & Secaucus NY4 & 550.2 & 1\,835.3 & 1\,835.9 & 2\,683.3 \\
Markham & Carteret & 553.1 & 1\,844.9 & 1\,845.4 & 2\,697.2 \\
\bottomrule
\end{tabular}
\caption{Distances and one-way \omterm{def:nw:the-north-american-map:floor}{latency floors} between the sites of \cref{dat:nw:the-north-american-map:sites}: WGS-84 geodesics (Vincenty's
inverse method) and the time light needs along them in vacuum, in air ($n = 1.0003$) and in standard single-mode fibre ($n_g = 1.4620$).
Data: \texttt{nw\_map.pair\_rows()}.}\label{tab:nw:the-north-american-map:pairs}
\end{table}

\section{The Chicago-area futures site}

CME Group's 10-K calls the Aurora building its largest data centre, leased from CyrusOne under a lease that expires in 2031, and describes a
private Google Cloud region and co-location facility that CME and Google are building in Aurora to host its markets; chapter~16 returns to
what a matching engine in a cloud region changes. Aurora has not always been the site. Laughlin, Aguirre and Grundfest, reading the
exchanges' own timestamps, place CME's move of its matching engine from 350 Cermak, in the city, to Aurora in 2010, and the move of CME's
\omterm{def:nw:colocation-products-and-how-they-are-sold:colo}{colocation} to Aurora on 29 January 2012; until then a firm colocated at Cermak with a microwave link that started in Aurora paid a metro round
trip between the two. The distance between them is \qty{52.3}{\kilo\metre}, \qty{255}{\micro\second} of straight fibre each way, so the round
trip in fibre was at least half a millisecond, the figure their paper puts on it.

The Chicago side is thus one building for futures and a second site in the city, which remains Cboe's secondary data centre for all its U.S.
platforms and a point of presence for networks. The continent's most-watched race runs from that single building to three: equity index
futures prices formed in Chicago lead the cash prices formed in New Jersey, as the same paper recalls, and every firm that trades on the lead
needs the fastest path from Aurora to each corner of the triangle.

\begin{omfigure}
\begin{tikzpicture}
\begin{axis}[axis equal image,width=12.5cm,xmin=-850,xmax=1120,ymin=-230,ymax=320,xlabel={km east of 81\textdegree{} W},
  ylabel={km north of 41.8\textdegree{} N},tick label style={font=\footnotesize,/pgf/number format/1000 sep={\,}},label style={font=\small},grid=major,grid style={gray!20},
  table/col sep=comma,unbounded coords=jump]
\addplot[omDef,thick,no markers] table[x=x,y=y]{figdata/networks/10-the-north-american-map/na_links.csv};
\addplot[only marks,mark=*,mark size=2.4pt,omThm,visualization depends on={value \thisrow{anchor}\as\anc},nodes near coords,
  point meta=explicit symbolic,every node near coord/.append style={font=\footnotesize,anchor=\anc,black}]
  table[x=x,y=y,meta=name]{figdata/networks/10-the-north-american-map/na_sites.csv};
\addplot[only marks,mark=none,nodes near coords,point meta=explicit symbolic,
  every node near coord/.append style={font=\scriptsize,fill=white,inner sep=1pt,anchor=center,text=omMeth}]
  table[x=x,y=y,meta=text]{figdata/networks/10-the-north-american-map/na_labels.csv};
\end{axis}
\end{tikzpicture}

\omcaption{Chicago, Toronto and New Jersey from computed coordinates, each link labelled with its \omterm{def:nw:the-north-american-map:geodesic}{geodesic distance} and light-time in vacuum.
At this scale the New Jersey triangle of \cref{fig:nw:the-north-american-map:nj} is two dots (Mahwah above, Carteret below). A local
equirectangular projection centred on 41.8\textdegree{}~N, 81\textdegree{}~W; over 1\,500~km it stretches shapes by a few per cent, which a
schematic tolerates and a distance computation does not (distances are computed on the ellipsoid). Data: \texttt{fig\_map.py} on
\texttt{firm.geomap}.}\label{fig:nw:the-north-american-map:na}
\end{omfigure}

\section{Toronto and the rest of the continent}

TMX runs its own data centre in Markham, north of Toronto: its pages describe a suburban building operated for capital-markets clients,
carrier- and vendor-neutral, with metro millimetre-wave and cross-border microwave services on site. ICE's page speaks of the inter-listed equities
and exchange-traded funds that trade in both Toronto and New York, and sells a wireless service between them: TSX data delivered to Mahwah, to Secaucus NY4 and to Carteret, U.S. data delivered to Markham, with private bandwidth of 1 to
10~megabits a second, and publishes three latencies per route: radio to radio, rack to rack, and the figure its service-level agreement
guarantees. From Markham to Mahwah they are 1.88, 1.89 and \qty{1.90}{\milli\second}.

The page does not say whether these are one-way or round-trip times, and it need not: the geodesic answers. From Markham to Mahwah light in
vacuum needs \qty{1.747}{\milli\second} one way, so a round trip cannot take less than \qty{3.49}{\milli\second}, and \qty{1.89}{\milli\second}
can only be one way, 8\% above the floor in air. The rack-to-rack figure adds \qty{10}{\micro\second} to Mahwah and \qty{20}{\micro\second} to
Carteret over radio to radio: the short connections at each end, from the building's roof or tower to the customer's rack.

\begin{remark}[The rest of the continent]\label{rem:nw:the-north-american-map:rest}
The method of this chapter extends to any site whose building can be sourced: the options exchanges, the other futures venues, the markets
of Mexico. \texttt{firm.venuesites} (chapter~11) holds the registry from venue to site; chapters~22 to~26 fill it for the venues
each asset class needs. What does not extend is the certainty: for some buildings a filing or an operator's page names the street address, for
others only a directory does, and the site table must say which.
\end{remark}

\section{Distances, light-times and route factors}

\begin{definition}[Latency floor, group index]\label{def:nw:the-north-american-map:floor}
The \emph{latency floor}\index{latency floor} between two sites is the one-way time light needs along the geodesic between them, $d/c_0$ in
vacuum, with $c_0 = \qty{299792458}{\metre\per\second}$; in a medium it is $n_g d / c_0$. The \emph{group index}\index{group index} $n_g$ of a
medium is the ratio of $c_0$ to the speed at which a pulse, and so information, travels through it: about 1.462 in standard single-mode fibre
at \qty{1550}{\nano\metre}, about 1.0003 in air near the ground.
\end{definition}

\begin{definition}[Route factor]\label{def:nw:the-north-american-map:factor}
The \emph{route factor}\index{route factor} of a route is its published or measured one-way latency divided by the \omterm{def:nw:the-north-american-map:floor}{latency floor}, in the route's
own medium, between its end points. A route factor of 1 is a route along the geodesic with no equipment; everything else is above 1.
\end{definition}

\begin{proposition}[What the floors are worth]\label{prop:nw:the-north-american-map:floors}
Per \qty{1000}{\kilo\metre} of geodesic the \omterm{def:nw:the-north-american-map:floor}{latency floor} is \qty{3.336}{\milli\second} in vacuum and in air to within
\qty{1}{\micro\second}, and \qty{4.877}{\milli\second} in fibre: along the same path, glass is 46\% slower than air.
A route of factor $\rho$ in a medium of \omterm{def:nw:the-north-american-map:floor}{group index} $n_g$ spends $(\rho - 1)$ times its floor on everything but the geodesic, the equivalent of
$(\rho - 1)\,d$ kilometres of extra path in that medium.
\end{proposition}

\begin{proof}
$10^3/c_0$ is \qty{3.3356}{\micro\second} per kilometre; times $1.0003$ it grows by \qty{1.0}{\micro\second} per \qty{1000}{\kilo\metre}; times
1.4620 it is \qty{4.877}{\micro\second\per\kilo\metre}, 46.2\% more. A time $\rho\, n_g d / c_0$ exceeds the floor by $(\rho-1)\,n_g d/c_0$,
which is the time the same medium needs for $(\rho-1)\,d$ more kilometres.
\end{proof}

The air's own index barely matters: the Ciddor equation, as NIST tabulates it, gives 1.00027 for dry air at \qty{20}{\celsius} and
sea-level pressure and 1.00030 for warm, humid, compressed air, a spread that moves the Aurora--Carteret floor by \qty{0.1}{\micro\second}.
What matters is glass against air, \qty{1.82}{\milli\second} on that route, and how far the route strays from the geodesic.

\omcode{code/firm/geomap/firm_geomap.py}{83}{89}{Latency floors and route factors: the medium's index times the geodesic over the speed of
light, and a published latency over the floor.}\label{lst:nw:the-north-american-map:floor}

\Cref{tab:nw:the-north-american-map:routes} applies the definitions to the published routes in the ledger. Laughlin, Aguirre and Grundfest's
2010 fibre, between 350 Cermak and Carteret, ran at a quoted \qty{6.65}{\milli\second}, later \qty{6.55}{\milli\second}: a \omterm{def:nw:the-north-american-map:factor}{route factor} of 1.21,
then 1.19, over its fibre floor, or \qty{1.1}{\milli\second} of extra time, the equivalent of 234 kilometres of glass. Some of it is path (a
cable follows roads and rights of way), some is equipment (amplifiers, regenerators, the switches at each end). The same authors estimated the
microwave networks licensed in 2011 and 2012 at 4.2 to \qty{5.2}{\milli\second}; the network that reached Carteret in \qty{3.982}{\milli\second}
in 2016 had a \omterm{def:nw:the-north-american-map:factor}{route factor} of 1.011, \qty{42}{\micro\second} above the floor in air, or about 12.5~kilometres of detour and equipment over
1\,181. Chapter~13 takes the fibre routes apart and chapter~14 the towers.

\begin{table}[ht]
\centering\small
\begin{tabular}{lllrrrr}
\toprule
Route & From & To & Published & Floor & Route & Excess \\
 & & & (ms) & (ms) & factor & (\si{\micro\second}) \\
\midrule
Microwave, 2016 & Aurora & Carteret & 3.982 & 3.940 & 1.011 & 42 \\
Microwave, 2016 & Aurora & Mahwah & 3.986 & 3.934 & 1.013 & 52 \\
Wireless, rack to rack & Markham & Mahwah & 1.890 & 1.748 & 1.082 & 142 \\
Wireless, rack to rack & Markham & NY4 & 1.990 & 1.836 & 1.084 & 154 \\
Wireless, rack to rack & Markham & Carteret & 2.050 & 1.845 & 1.111 & 205 \\
Fibre, 2010 & 350 Cermak & Carteret & 6.650 & 5.507 & 1.208 & 1\,143 \\
Fibre, later & 350 Cermak & Carteret & 6.550 & 5.507 & 1.189 & 1\,043 \\
\bottomrule
\end{tabular}
\caption{Published one-way latencies against the floor in each route's medium (air for microwave and wireless, fibre for fibre). Sources:
Quincy Data's 2016 release for the microwave routes, ICE's Toronto wireless page for the Markham routes, Laughlin, Aguirre and Grundfest
(2014) for the fibre. Data: \texttt{nw\_map.route\_rows()}.}\label{tab:nw:the-north-american-map:routes}
\end{table}

\omcode{code/networks/10-the-north-american-map/python/nw_map.py}{63}{75}{Each published route against its floor: the route factor, the excess
time and the extra path in the route's own medium that the excess is worth.}\label{lst:nw:the-north-american-map:routes}

\begin{omfigure}
\begin{tikzpicture}
\begin{axis}[width=11.5cm,height=6.8cm,xmin=0,xmax=1250,ymin=0,ymax=7.2,xlabel={geodesic distance (km)},
  ylabel={one-way latency (ms)},tick label style={font=\footnotesize,/pgf/number format/1000 sep={\,}},label style={font=\small},grid=major,grid style={gray!25},
  legend pos=south east,legend style={font=\footnotesize},table/col sep=comma]
\addplot[omMeth,thick,dashed,no markers] table[x=km,y=fibre_ms]{figdata/networks/10-the-north-american-map/floors.csv};
\addplot[black,thick,no markers] table[x=km,y=vacuum_ms]{figdata/networks/10-the-north-american-map/floors.csv};
\addplot[only marks,mark=*,mark size=2.2pt,omThm] table[x=km,y=published_ms]{figdata/networks/10-the-north-american-map/routes_air.csv};
\addplot[only marks,mark=square*,mark size=2.2pt,omDef] table[x=km,y=published_ms]{figdata/networks/10-the-north-american-map/routes_fibre.csv};
\legend{floor in fibre,floor in vacuum,published: air,published: fibre}
\end{axis}
\end{tikzpicture}

\omcaption{Published one-way route latencies against the \omterm{def:nw:the-north-american-map:geodesic}{geodesic distance} between their end points, with the floors in vacuum (the floor in air
is indistinguishable at this scale) and in fibre. The wireless routes sit just above the vacuum line; the 2010 fibre sits above even the fibre
floor. Data: \texttt{fig\_map.py}, routes of \cref{tab:nw:the-north-american-map:routes}.}\label{fig:nw:the-north-american-map:routes}
\end{omfigure}

\begin{remark}[Why the ellipsoid]\label{rem:nw:the-north-american-map:sphere}
The spherical (haversine) formula with the Earth's mean radius puts Aurora and Carteret \qty{2.97}{\kilo\metre} closer than the ellipsoid does, an
error of 0.25\% and \qty{9.9}{\micro\second} of light in air: more than a quarter of the gap between the 2016 microwave route and its floor.
Between Mahwah and Carteret the sphere errs by 0.12\% the other way. \omterm{def:nw:the-north-american-map:factor}{Route factors} of 1.01 cannot be measured with a yardstick that is wrong by
0.25\%, which is why \texttt{firm.geomap} computes Vincenty's geodesic on WGS-84 and tests it against Karney's published set of geodesics.
\end{remark}

\begin{omfigure}
\begin{tikzpicture}[font=\footnotesize,>=Stealth]
\draw[thick,gray!60,dashed] (0,0) -- (11,0);
\node[below,gray] at (5.5,-0.1) {geodesic: the latency floor};
\draw[thick,omDef] (0,0) -- (0.6,0.9) -- (2.4,1.2) -- (4.1,0.95) -- (5.9,1.35) -- (7.7,1.15) -- (9.3,1.0) -- (10.4,0.9) -- (11,0);
\foreach \x/\y in {0.6/0.9,2.4/1.2,4.1/0.95,5.9/1.35,7.7/1.15,9.3/1.0,10.4/0.9}{\draw[omThm,fill=omThm!30] (\x-0.09,\y-0.09) rectangle (\x+0.09,\y+0.09);}
\node[draw,fill=white,rounded corners=2pt,inner sep=2pt,align=center] at (0,-0.9) {rack in\\ data centre A};
\node[draw,fill=white,rounded corners=2pt,inner sep=2pt,align=center] at (11,-0.9) {rack in\\ data centre B};
\draw[<->,omMeth] (0.1,0.12) -- node[left,align=right,font=\scriptsize]{tail} (0.52,0.76);
\draw[<->,omMeth] (10.9,0.12) -- node[right,align=left,font=\scriptsize]{tail} (10.48,0.76);
\node[omThm,font=\scriptsize] at (5.9,1.8) {towers, amplifiers or regenerators};
\node[omDef,font=\scriptsize] at (5.5,0.45) {path: detours around terrain and rights of way};
\end{tikzpicture}

\omcaption{What a \omterm{def:nw:the-north-american-map:factor}{route factor} contains, schematically: the detour of the path from the geodesic, the equipment along it (towers, amplifiers,
regenerators) and the tails at each end, from the building's roof or \omterm{def:nw:colocation-products-and-how-they-are-sold:mmr}{meet-me room} to the customer's rack. A published figure that is ``radio to
radio'' leaves the tails out; ``rack to rack'' includes them.}\label{fig:nw:the-north-american-map:parts}
\end{omfigure}

\section{Method: a site table you can defend}

\begin{method}[Building the site table]\label{met:nw:the-north-american-map:table}
\begin{enumerate}
\item \textbf{Source the building, not the city.} Take the site from the venue's own documents (a connectivity manual, a filing, an annual report)
  and the street address from the venue or the data-centre operator. A directory is a lead, not a source; if it is all there is, flag the site.
\item \textbf{Geocode with a cited geocoder.} Record the geocoder, the query and the date; store latitude and longitude with seven decimals
  (about a centimetre), which is more than the address justifies but keeps the computation reproducible.
\item \textbf{Cross-check independently.} Compare with an independent published position of the same building; a disagreement of more than a
  few hundred metres means one of them is wrong.
\item \textbf{Compute, never type.} Distances, floors and map coordinates come from the table through code; the table's rows carry their ledger
  ids so that every number on a map can be traced to a source.
\item \textbf{Date it.} Venues move (CME in 2010 and 2012, as above; Euronext in 2022, in chapter~11); re-verify the table before every use
  of it that matters.
\end{enumerate}
\end{method}

How much does a site error cost? A kilometre is \qty{3.3}{\micro\second} of light in vacuum, which is the size of the gaps that the fastest
routes compete over. The site table of this chapter meets the method's third step: Laughlin, Aguirre and Grundfest give Aurora at about
41.80\textdegree{}~N, 88.24\textdegree{}~W and Carteret at 40.58\textdegree{}~N, 74.25\textdegree{}~W, which agree with the geocoded coordinates to
within a hundredth of a degree (under a kilometre). Their distance, \qty{1179}{\kilo\metre}, lies between the spherical (\qty{1177.2}{\kilo\metre}) and the ellipsoidal
(\qty{1180.1}{\kilo\metre}) distances between their rounded positions, and within \qty{2}{\kilo\metre} of ours; the paper does not say which
Earth it used, and its light-time of \qty{3.93}{\milli\second} agrees with ours to \qty{6}{\micro\second}. Markham is the weak row: its
building is sourced to TMX and ICE, its street address only to a directory, and a Markham error of a few kilometres would move its floors by
some \qty{10}{\micro\second}, which the \omterm{def:nw:the-north-american-map:factor}{route factors} of \cref{tab:nw:the-north-american-map:routes} would absorb without anyone noticing.

\omcode{code/firm/geomap/firm_geomap.py}{43}{73}{Vincenty's inverse method on the WGS-84 ellipsoid, with a convergence guard: the iteration on the
longitude difference on the auxiliary sphere, then the series for the distance.}\label{lst:nw:the-north-american-map:vincenty}

\section{Tutorial: the map from computed coordinates}\label{tut:nw:the-north-american-map:tut}

\begin{tutorial}
\textbf{Goal.} Load the site table, compute every distance and floor, compare published routes with their floors, and draw the maps.
\textbf{End state:} \cref{fig:nw:the-north-american-map:nj,fig:nw:the-north-american-map:na,fig:nw:the-north-american-map:routes} and
\cref{tab:nw:the-north-american-map:pairs,tab:nw:the-north-american-map:routes}.
\begin{tutsteps}
\item \textbf{The table.} \texttt{data/networks/sites.csv} holds seven sites, each with the ledger rows that source its building, address and
  coordinates; \texttt{firm\_geomap.load\_sites} reads it.
\item \textbf{The yardstick.} Run \texttt{firm.geomap}'s tests: Vincenty's distance
  (\cref{lst:nw:the-north-american-map:vincenty}) against 249 geodesics of Karney's test set, in Python and in C++20, to better than a tenth of a
  millimetre.
\item \textbf{The floors.} \texttt{nw\_map.pair\_rows()} computes \cref{tab:nw:the-north-american-map:pairs}; \texttt{route\_rows()}
  (\cref{lst:nw:the-north-american-map:routes}) computes \cref{tab:nw:the-north-american-map:routes}.
\item \textbf{The maps.} \texttt{python fig\_map.py} projects each region's sites and writes the site, link and label files that the TikZ
  maps read; no coordinate is typed in the figure.
\end{tutsteps}
\textbf{What to change next.} Add Cboe's secondary site's points of presence, or another venue from its own manual, with its ledger row, and
recompute; replace the fibre index by that of a hollow-core fibre (chapter~13) and see what happens to the Chicago floor.
\end{tutorial}

\section{Build: the geodesic map}\label{bld:nw:the-north-american-map:geomap}

\begin{build}
\bfield{Purpose} Sites as sourced data, distances on the ellipsoid, \omterm{def:nw:the-north-american-map:floor}{latency floors} and \omterm{def:nw:the-north-american-map:factor}{route factors}, and coordinates for maps: the base of
every map and route of Part~III, and of chapter~29's plan.
\bfield{Interface} \texttt{firm\_geomap}: \texttt{Site(id, name, operator, lat, lon, source, venues)}, \texttt{geodesic\_m},
\texttt{haversine\_m}, \texttt{floor\_us(d\_m, medium, n\_g)}, \texttt{route\_factor(published\_us, d\_m, medium)}, \texttt{matrix},
\texttt{project(sites, lat0, lon0)}, \texttt{load\_sites}, \texttt{write\_sites}; a C++20 twin of the geodesic in \texttt{cpp/firm\_geomap.hpp}.
\bfield{Rules} Every site names its source rows; the geodesic raises rather than returning a wrong distance when the iteration fails (nearly
antipodal points); floors are one-way; a projection is for drawing only, never for a distance.
\bfield{Acceptance tests} \texttt{code/firm/geomap/tests/}: 249 geodesics of Karney's CC0 test set to \qty{0.1}{\milli\metre} in Python and
C++20, the floors by hand, the sphere against the ellipsoid, and the site table's sources present.
\bfield{Stretch} Karney's algorithm, which converges everywhere; a great-circle path as a list of points, for drawing a route on a map.
\end{build}

\begin{omsources}
\item G.~Laughlin, A.~Aguirre and J.~Grundfest, ``Information transmission between financial markets in Chicago and New York'', \emph{Financial
Review} 49(2) (2014).
\item T.~Vincenty, ``Direct and inverse solutions of geodesics on the ellipsoid with application of nested equations'', \emph{Survey Review}
23(176) (1975); C.~F.~F.~Karney, ``Algorithms for geodesics'', \emph{Journal of Geodesy} 87 (2013), and the test set for geodesics (Zenodo, CC0).
\item ICE, \emph{Colocation Technical Specifications} (Mahwah) and the Toronto wireless page; Cboe, \emph{Titanium U.S. Equities/Options
Connectivity Manual}; CME Group, Form 10-K 2024; SEC Release 34-101267 (Nasdaq); Quincy Data, release of 13 May 2016.
\item BIPM, the SI definition of the metre; NIST, refractive index of air (Ciddor equation); site coordinates from OpenStreetMap Nominatim.
\end{omsources}

\section{Exercises}

\begin{exercise}[$\star$]\label{exo:nw:the-north-american-map:1}
How long does light need to cover \qty{1000}{\kilo\metre} in vacuum and in standard fibre? What is the difference over Aurora--Carteret?
\end{exercise}

\begin{exercise}[$\star$]\label{exo:nw:the-north-american-map:2}
A vendor publishes a latency of \qty{2.05}{\milli\second} between Markham and Carteret without saying whether it is one-way. Show which it is.
\end{exercise}

\begin{exercise}[$\star$]\label{exo:nw:the-north-american-map:3}
The spherical formula puts Aurora and Carteret \qty{2.97}{\kilo\metre} too close. How many microseconds in air is that, and why does it matter?
\end{exercise}

\begin{exercise}[$\star\star$]\label{exo:nw:the-north-american-map:4}
The 2016 microwave route from Aurora to Mahwah was published at \qty{3.986}{\milli\second}. Compute its \omterm{def:nw:the-north-american-map:factor}{route factor} and the extra path it is
worth in air.
\end{exercise}

\begin{exercise}[$\star\star$]\label{exo:nw:the-north-american-map:5}
A router in Carteret sends one piece of an order to Secaucus NY5 and one to Mahwah over straight fibre. By how much should synchronised routing
delay the NY5 piece, and why is the real figure different?
\end{exercise}

\begin{exercise}[$\star\star$]\label{exo:nw:the-north-american-map:6}
Which row of the site table do you trust least, and what would a 3-kilometre error in it do to the floors that use it?
\end{exercise}

\begin{exercise}[$\star\star\star$]\label{exo:nw:the-north-american-map:7}
\emph{Coding.} Round Aurora's and Carteret's coordinates to two decimals, as Laughlin, Aguirre and Grundfest print them, and compute the geodesic
with \texttt{firm.geomap}. How far is it from theirs and from the unrounded one?
\end{exercise}

\begin{exercise}[$\star\star\star$]\label{exo:nw:the-north-american-map:8}
\emph{Find the flaw.} ``Our fibre route between Chicago and New Jersey has a \omterm{def:nw:the-north-american-map:factor}{route factor} of 1.19, so a microwave network on the same path would
be 19\% faster.''
\end{exercise}

\section{Problem: Carteret to Aurora}

\begin{problem}[{Weekend problem --- the floors between the Nasdaq and CME data centres, and a fibre route against them}]\label{pb:nw:the-north-american-map:1}
A firm trades index futures in Aurora against stocks in Carteret and wants to know how fast the link between them could ever be, and how far
from that limit the published routes have been.

\medskip\noindent\textbf{Part I --- The distance.}
\begin{enumerate}
\item From the site table, what is the \omterm{def:nw:the-north-american-map:geodesic}{geodesic distance} between Aurora and Carteret?
\item What does the spherical formula give, and what is its error in kilometres and per cent?
\item How much shorter is the straight chord through the Earth?
\item Why is the geodesic, not the chord, the right yardstick for a trading route?
\end{enumerate}

\medskip\noindent\textbf{Part II --- The floors.}
\begin{enumerate}[resume]
\item What is the one-way \omterm{def:nw:the-north-american-map:floor}{latency floor} in vacuum?
\item In air, taking $n = 1.0003$, and how much does the choice between 1.00027 and 1.0003 matter?
\item In standard fibre along the geodesic?
\item What is the round-trip floor in fibre, and in vacuum?
\end{enumerate}

\medskip\noindent\textbf{Part III --- A fibre route against its floor.}
\begin{enumerate}[resume]
\item The 2010 fibre from 350 Cermak to Carteret was quoted at \qty{6.65}{\milli\second} one way. What is its floor, and its \omterm{def:nw:the-north-american-map:factor}{route factor}?
\item How much excess time is that, and how many kilometres of glass is it worth?
\item From Aurora, a firm reached that fibre over the Cermak--Aurora metro. If the metro leg had the same \omterm{def:nw:the-north-american-map:factor}{route factor}, what would the time be
  from Aurora to Carteret, and its \omterm{def:nw:the-north-american-map:factor}{route factor} against the Aurora--Carteret floor in fibre?
\item The route was later quoted at \qty{6.55}{\milli\second}. What changed in the \omterm{def:nw:the-north-american-map:factor}{route factor}?
\end{enumerate}

\medskip\noindent\textbf{Part IV --- The verdict.}
\begin{enumerate}[resume]
\item State the \emph{named result}: the three floors between Carteret and Aurora and the \omterm{def:nw:the-north-american-map:factor}{route factor} of the 2010 fibre.
\item The 2016 microwave route reached Carteret in \qty{3.982}{\milli\second}. By how much did it beat the fibre floor?
\item What is left between it and the floor in air?
\item Could any fibre route ever beat it? What kind of fibre would it take (chapter~13)?
\item Which site error would move the named result by more than \qty{1}{\micro\second}?
\item What does the move of CME's matching engine from Cermak to Aurora do to the Cermak-based fibre's value?
\item Why do firms publish, and buy, rack-to-rack latencies rather than floors?
\item In one sentence: what does the geodesic tell a firm that a vendor's latency does not?
\end{enumerate}
\end{problem}

\section{Interview questions}

\begin{interviewq}[$\star$ \iqroles{developer, trader}]\label{iq:nw:the-north-american-map:1}
Where are the main U.S. equity, options and futures matching engines, and roughly how far apart?
\end{interviewq}

\begin{interviewq}[$\star$ \iqroles{developer}]\label{iq:nw:the-north-american-map:2}
What is the minimum one-way latency between Chicago and New Jersey, in vacuum and in fibre? Show the arithmetic.
\end{interviewq}

\begin{interviewq}[$\star\star$ \iqroles{developer, researcher}]\label{iq:nw:the-north-american-map:3}
Why is microwave faster than fibre over the same path, and why is fibre still used?
\end{interviewq}

\begin{interviewq}[$\star\star$ \iqroles{developer}]\label{iq:nw:the-north-american-map:4}
How would you compute the distance between two data centres, and how accurate does it need to be?
\end{interviewq}

\begin{interviewq}[$\star\star$ \iqroles{trader, researcher}]\label{iq:nw:the-north-american-map:5}
Why does a trader in Carteret see Nasdaq's quotes before the consolidated feed shows them, and how much of the gap is physics?
\end{interviewq}

\begin{interviewq}[$\star\star\star$ \iqroles{developer, trader}]\label{iq:nw:the-north-american-map:6}
A vendor offers a Chicago--New Jersey link at \qty{3.99}{\milli\second}. How do you judge the claim before you buy?
\end{interviewq}
